\section{--nmagic, or the Linker's Idea of a Segment}
\label{sec:build-nmagic}

This section documents the least intuitive flag in the whole chain, because without it the image is
built, verified, and written to the wrong place.

\subsection{What the Linker Does by Default}

ELF loadable segments carry two addresses: a virtual address (\codeinline{VMA}, where the code expects
to run) and a physical address (\codeinline{LMA}, where the bytes must land in memory). For flash that
is always 1:1 — you cannot relocate code that runs from flash.

A linker aiming at an operating system still aligns segments to page boundaries, and to do that it
gives the segment an \emph{offset} that no longer matches its load address. An OS loader compensates;
a bare-metal programmer does not, and writes the file's bytes at the segment's file offset while the
code expects them at the segment's address.

\subsection{The Symptom, Measured}

\begin{minted}{text}
# with a default link: the loadable segment's file offset is 0x1000, its address is 0x08006000
probe-rs download image.elf        # writes bytes at the wrong offset
probe-rs verify image.elf          # verifies the same wrong offsets: it agrees with itself
                                   # the board does not run
\end{minted}

This is the failure that teaches the difference between ``verified'' and ``correct''. The programmer
compares the flash against the file it just wrote. If the file was written at the wrong offset by a
consistent rule, both sides agree and the check passes.

\subsection{The Fix, and Why It Is One Flag}

\begin{minted}{text}
# in .cargo/config.toml, or as a link argument
rustflags = ["-C", "link-arg=--nmagic"]
\end{minted}

\codeinline{--nmagic} tells the linker not to align segments and not to pad them — the image is flat,
VMA equals LMA for every loadable segment, and the programmer's naive write is correct. Every bare
metal Rust project needs it and most tutorials omit it, because their toolchain's runner happens to
compensate.

\subsection{Checking the Invariant Instead of Trusting the Flag}

The project's inspection script prints exactly what matters, so the flag is verified rather than
believed:

\begin{minted}{bash}
python3 scripts/elf-load-range.py "$ELF"
# for each PT_LOAD: paddr == vaddr, and load_min >= 0x08006000
#   seg paddr=0x08006158 vaddr=0x08006158 filesz=5400
#   seg paddr=0x08007670 vaddr=0x08007670 filesz=108
#   seg paddr=0x080076DC vaddr=0x20000000 filesz=4     <- the .data initialiser
\end{minted}

Three lines, and two facts read directly off them: every code segment has \codeinline{paddr == vaddr}
(the flag worked), and \codeinline{load\_min} is above the bootloader (the memory map is right). The
third segment is the initialiser for \codeinline{.data}, which correctly has a flash address and a RAM
address — that one is supposed to differ.
